% 87-01 ④(87쪽) — 보기 그래프: $x=-2$ 에서 기울기 0, $x=1$ 에서 극대인 위로 볼록한 사차함수 개형. 그림 자체가 보기라 TikZ 로 다시 그림.
% 라벨은 원문 그대로 -2, O, 1, x, y, y=f(x) 와 점선 두 줄만.
\begin{tikzpicture}[baseline=(current bounding box.center), x=0.40cm, y=0.40cm, line width=0.5pt]
  \path (-3.6,-2.2) rectangle (3.45,3.6);
  \draw[->, >=stealth, line width=0.7pt] (-3.05,0) -- (2.7,0) node[below=1pt, font=\footnotesize] {$x$};
  \draw[->, >=stealth, line width=0.7pt] (0,-1.9) -- (0,3.4) node[left=1pt, font=\footnotesize] {$y$};
  \draw[densely dotted, line width=0.4pt] (-2,0) -- (-2,0.15);
  \draw[densely dotted, line width=0.4pt] (1,0) -- (1,2.62);
  \draw[line width=0.6pt, smooth] plot coordinates {(-2.75,-1.6) (-2.65,-1.1) (-2.5,-0.45) (-2.35,-0.05) (-2.2,0.1) (-2,0.15) (-1.8,0.2) (-1.6,0.3) (-1.35,0.5) (-1.1,0.75) (-0.8,1.15) (-0.5,1.55) (-0.2,1.95) (0.1,2.25) (0.4,2.48) (0.7,2.6) (1,2.62) (1.2,2.5) (1.4,2.2) (1.6,1.6) (1.75,1.0) (1.9,0.2) (2,-0.5) (2.1,-1.3) (2.15,-1.8)};
  \node[below=1pt, font=\footnotesize] at (-2,0) {$-2$};
  \node[below left=-2pt and -3pt, font=\footnotesize] at (0,0) {$\pt{O}$};
  \node[below=1pt, font=\footnotesize] at (1,0) {$1$};
  \node[anchor=west, font=\footnotesize] at (0.28,3.2) {$y=f(x)$};
\end{tikzpicture}
